\section{A tower obstruction in dimension two}
\label{sec:tower}

We first isolate the obstruction to a finite two-dimensional model.
Its input is the vanishing of second homology with a rank-one local
system.  Its conclusion restricts the maps into the complex from finite
complexes with vanishing first rational Betti number.  The case of a
Moore complex will then show that the corresponding character cannot
detect a finite-order element.

A \emph{rank-one complex local system} $L$ on a connected CW complex $T$
assigns a one-dimensional complex vector space to each point, with
parallel transport along paths that depends only on homotopy relative to
the endpoints.  After choosing a basepoint and a basis of its fiber,
its monodromy is a character
\[
  \rho:\pi_1(T)\longrightarrow\C^\times.
\]
We use ordinary homology with these coefficients.  In module notation,
write $\C_\rho$ for the right $\Z[\pi_1(T)]$-module with
$v\cdot g=\rho(g)v$.  If the universal-cover cellular chains carry the
left deck action, then
\[
  C_*(T;L)
  =\C_\rho\otimes_{\Z[\pi_1(T)]}C_*(\widetilde T;\Z).
\]
Thus cellular chains have finite support, even when $T$ is infinite.
Equivalently, $C_j(T;L)$ is the direct sum of one coefficient fiber for
each oriented $j$-cell, with boundary maps defined by parallel transport.

Two elementary consequences will be used throughout the proof.  First,
if $A\subseteq T$ is a subcomplex and $\dim T\leq2$, then
\begin{equation}
  H_2(A;L|_A)\longrightarrow H_2(T;L)
  \quad\text{is injective}.
  \label{eq:subcomplex-injection}
\end{equation}
Indeed, the cellular chain inclusion is injective, and in dimension at
most two the second homology is the kernel of the second boundary.
This argument requires no injectivity of $\pi_1(A)\to\pi_1(T)$.
Second, on a connected complex the zeroth homology is the space of
coinvariants,
\[
  H_0(T;L)
  =\C/\langle (\rho(g)-1)v:g\in\pi_1(T),\ v\in\C\rangle.
\]
Consequently, nontrivial monodromy implies $H_0(T;L)=0$.

The next lemma ensures that vanishing of second homology also survives
the covering operation needed in the tower.

\begin{lemma}[Infinite cyclic covers]
\label{lem:cyclic-cover}
Let $A$ be a finite connected CW complex of dimension at most two, let
$L$ be a rank-one complex local system on $A$, and let
$p:\widehat A\to A$ be a connected infinite cyclic cover.  If
$H_2(A;L)=0$, then
\[
  H_2(\widehat A;p^*L)=0.
\]
\end{lemma}

\begin{proof}
Choose a generator $\tau$ of the deck group and put
$R=\C[z,z^{-1}]$.  There is a canonical deck action on the pulled-back
local system: in the fiber description of $p^*L$, it is
\begin{equation}
  (\widehat x,v)\longmapsto(\tau\widehat x,v),
  \qquad v\in L_{p(\widehat x)}.
  \label{eq:canonical-deck-action}
\end{equation}
The equality $p\tau=p$ makes this action compatible with parallel
transport.  Let $z$ act by this deck transformation.  Choosing one lift
of each cell of $A$ and a basis in its coefficient fiber identifies
$C_j(\widehat A;p^*L)$ with a finite free $R$-module, of rank equal to
the number of $j$-cells of $A$.

Taking coinvariants of this deck action identifies translated cell
lifts and their coefficient fibers.  It therefore gives an isomorphism
of chain complexes
\begin{equation}
  C_*(\widehat A;p^*L)\otimes_R R/(z-1)
  \cong C_*(A;L).
  \label{eq:cyclic-specialization}
\end{equation}
In the chosen bases, let $D(z)$ be the matrix of the second boundary
upstairs, and let $m$ be its number of columns.  The matrix $D(1)$ is
the second boundary downstairs, so it is injective by the hypothesis.
If $m=0$, there is nothing to prove.  Otherwise, some $m\times m$ minor
of $D(1)$ is nonzero.  The corresponding minor of $D(z)$ is consequently
a nonzero Laurent polynomial.  Thus $D(z)$ has full column rank over
$\C(z)$ and is injective over $R$, since $R^m$ embeds in $\C(z)^m$.
There are no cells above dimension two, and hence
$H_2(\widehat A;p^*L)=\ker D(z)=0$.
\end{proof}

The following proof adapts the cyclic-tower method of
Howie~\cite{Howie81Pairs}, in which one alternates subcomplex inclusions
and infinite cyclic covers; see Louder--Wilton~\cite[Section~3]{LouderWilton17}.
Lemma~\ref{lem:cyclic-cover} supplies the property that persists along
the tower here: vanishing of twisted second homology.

\begin{theorem}[Tower obstruction]
\label{thm:tower}
Let $Y$ be a finite connected CW complex of dimension at most two, and
let $L$ be a rank-one complex local system on $Y$ with $H_2(Y;L)=0$.
There is no map $f:F\to Y$ from a finite connected simplicial complex
satisfying $b_1(F;\Q)=0$ for which $f^*L$ has nontrivial monodromy.
\end{theorem}

\begin{proof}
Suppose that such a map exists.  We will repeatedly lift it to an
infinite cyclic cover and replace that cover by the image of the lift.
At every step the image will contain strictly more simplices, although
the domain retains one fixed finite triangulation.  For a finite
simplicial complex $A$, write $s(A)$ for its number of nonempty simplices.

We first justify working simplicially without increasing the dimension
of $Y$.  Its one-skeleton is a finite graph.  Subdivide that graph to
make it simplicial and homotope each two-cell attaching map to an edge
path.  After further subdivision if necessary, each nonconstant
attaching path can be represented by a map from a polygonal circle
with at least three edges, with every edge mapping homeomorphically
onto a graph edge.  Replace the attachment by the mapping cylinder of this
map, capped by a disk along its free boundary circle.  This merely
inserts a collar into the attached disk.  The resulting cylinder has
one rectangle for each traversal of a graph edge.  Each rectangle is
embedded, with its two upper vertices in the separate domain circle
and its two lower vertices at distinct endpoints of the graph edge.
After compatible subdivisions of boundary edges, triangulate each
rectangle from its own interior vertex and each capping disk from a
new interior vertex.  This yields a finite simplicial complex even
when different rectangles traverse the same graph edge: their
interiors and upper edges remain distinct.  A constant attaching map
contributes a triangulated two-sphere joined to the graph at one
vertex.  Homotoping the attaching maps preserves the homotopy type,
so this construction gives a finite simplicial model of dimension at
most two.

Transfer $L$ and $f$ across this homotopy equivalence.  Homotopy
invariance preserves both $H_2(Y;L)=0$ and the nontrivial pullback
monodromy.  After a sufficiently fine subdivision of $F$, simplicial
approximation makes $f$ simplicial.  Fix this triangulation of $F$
for the rest of the proof.  Let $A_0=f(F)$, let $f_0:F\to A_0$ be the
corestriction, and let $L_0=L|_{A_0}$.  The simplicial image $A_0$ is
a finite connected subcomplex, and
\eqref{eq:subcomplex-injection} gives $H_2(A_0;L_0)=0$.

Here is the inductive step.  Suppose that $A_i$ is a finite connected
simplicial complex of dimension at most two, $f_i:F\to A_i$ is a
simplicial map onto its image $A_i$, and $L_i$ satisfies
\begin{equation}
  H_2(A_i;L_i)=0,
  \qquad f_i^*L_i\text{ has nontrivial monodromy}.
  \label{eq:tower-invariants}
\end{equation}
The second condition implies that $L_i$ itself has nontrivial
monodromy, so $H_0(A_i;L_i)=0$.  Since rank-one cellular chains have
the same dimensions as ordinary complex cellular chains, the Euler
characteristic gives
\[
  \chi(A_i)=-\dim_\C H_1(A_i;L_i)\leq0.
\]
On the other hand,
\[
  \chi(A_i)=1-b_1(A_i;\Q)+b_2(A_i;\Q).
\]
It follows that $b_1(A_i;\Q)\geq1$.  The finitely generated
abelianization of $\pi_1(A_i)$ therefore has a nonzero free part,
and there is an epimorphism
$\phi_i:\pi_1(A_i)\twoheadrightarrow\Z$.

Let $p_i:\widehat A_i\to A_i$ be the connected cover corresponding to
$\ker\phi_i$.  The composite $\phi_i(f_i)_*$ is zero: every
homomorphism $\pi_1(F)\to\Z$ factors through $H_1(F;\Z)$, which is
finite because $F$ is finite and $b_1(F;\Q)=0$.
Thus $f_i$ lifts to $\widetilde f_i:F\to\widehat A_i$.
Give the cover the lifted triangulation.  The lift is simplicial on
the already fixed triangulation of $F$, since on each domain simplex
it is the lift of a simplicial map into a single target simplex.
No new subdivision of $F$ is needed.

Set
\[
  A_{i+1}=\widetilde f_i(F),
  \qquad L_{i+1}=(p_i^*L_i)|_{A_{i+1}},
\]
and corestrict the lift to $f_{i+1}:F\to A_{i+1}$.
Lemma~\ref{lem:cyclic-cover}, followed by
\eqref{eq:subcomplex-injection}, gives
$H_2(A_{i+1};L_{i+1})=0$.  Moreover,
$f_{i+1}^*L_{i+1}\cong f_i^*L_i$, so the nontrivial monodromy
persists.  These are precisely the conditions
\eqref{eq:tower-invariants} needed to continue.  Figure~\ref{fig:tower}
shows one step.

\begin{figure}[tb]
  \centering
  \begin{tikzpicture}[
  >=Latex,
  every node/.style={font=\normalsize},
  every edge/.style={draw,->,semithick},
  baseline=(current bounding box.center)
]
  \node (F) at (0,0) {$F$};
  \node (next) at (3,1.1) {$A_{i+1}$};
  \node (cover) at (6.2,1.1) {$\widehat A_i$};
  \node (base) at (6.2,-1.1) {$A_i$};
  \path (F) edge node[above,sloped] {$f_{i+1}$} (next);
  \draw[semithick,{Hooks[right]}-Latex]
    (next) -- node[above] {inclusion} (cover);
  \path (cover) edge node[right] {$p_i$} (base);
  \node[anchor=west,font=\small,align=left] at (6.85,0)
    {infinite cyclic\\cover};
  \path (F) edge node[below,sloped] {$f_i$} (base);
  \node[font=\small] at (2.85,-1.85)
    {$s(A_i)<s(A_{i+1})\leq s(F)$};
\end{tikzpicture}

  \caption{One step of the tower.  The lifted map factors through its
  finite image $A_{i+1}\subseteq\widehat A_i$, and the diagram commutes.
  The domain $F$ keeps the same triangulation at every step; only the
  indicated subcomplex arrow is an inclusion.}
  \label{fig:tower}
\end{figure}

It remains to show that $s(A_i)$ increases strictly.  The restriction
$p_i|_{A_{i+1}}$ maps onto $A_i$ and maps each
simplex isomorphically onto a simplex.  Hence
$s(A_{i+1})\geq s(A_i)$.  If equality held, there would be exactly
one simplex above each simplex of $A_i$, compatibly with all faces.
The restriction would then be a simplicial isomorphism.  Its inverse,
followed by the inclusion into $\widehat A_i$, would give a section
of $p_i$.  Such a section would make $(p_i)_*$ surjective, contrary
to
\[
  \im(p_i)_*=\ker\phi_i\ne\pi_1(A_i).
\]
Thus $s(A_{i+1})>s(A_i)$.  But every simplex of $A_i$ is the image
of a simplex in the fixed triangulation of $F$, so
\[
  s(A_0)<s(A_1)<s(A_2)<\cdots\leq s(F).
\]
This is impossible, and the theorem follows.
\end{proof}

\Needspace{5\baselineskip}
The following consequence is the form used in the construction.

\begin{corollary}[Characters and torsion]
\label{cor:torsion}
Let $Y$ be a finite connected CW complex of dimension at most two.
If a character $\rho:\pi_1(Y)\to\C^\times$ satisfies
$H_2(Y;\C_\rho)=0$, then $\rho(g)=1$ for every finite-order
element $g\in\pi_1(Y)$.
\end{corollary}

\begin{proof}
Suppose that $g$ has finite order $n$ and $\rho(g)\ne1$; in
particular, $n\geq2$.  Let
\[
  M_n=S^1\cup_{\deg n}D^2
\]
be the Moore complex obtained by attaching a disk along a map of
degree $n$.  A loop in $Y$ representing $g$ extends to a map
$M_n\to Y$, because its $n$th power is nullhomotopic.  The cellular
boundary $C_2(M_n;\Z)\to C_1(M_n;\Z)$ is multiplication by $n$,
so $H_1(M_n;\Z)=\Z/n$ and $b_1(M_n;\Q)=0$.  The pullback
character sends the circle generator to $\rho(g)\ne1$.

Use the finite simplicial model for $M_n$ constructed in the proof of
Theorem~\ref{thm:tower}, and compose its homotopy equivalence to
$M_n$ with the map to $Y$.  This gives a finite connected simplicial
domain with vanishing first rational Betti number and nontrivial
pullback monodromy, contradicting that theorem.  For $n=2$ the
standard model $M_2$ is $\RP^2$, the case needed below.
\end{proof}
